%! AP Statistics — Unit 2, Lesson 2.1 — Exit Ticket KEY
\documentclass[10pt]{article}
\usepackage{apstats-article}
\usepackage{apstats-key}

%=============================================================================
\begin{document}

\pageheader{Unit 2, Lesson 2.1}{Exit Ticket --- Answer Key}
\begin{tabularx}{\linewidth}{X X X}
Name:\;\ans{Key} & Date:\; & Period:\;
\end{tabularx}

\vspace{0.22in}

\begin{tcolorbox}[colback=sky, colframe=navy, boxrule=0.9pt, arc=2mm,
    left=4mm, right=4mm, top=3mm, bottom=3mm]
\small\textit{A researcher surveys 30 high school students and finds that students
who drink more coffee (in cups per week) tend to have higher GPAs.}
\end{tcolorbox}

\vspace{0.18in}

\renewcommand{\arraystretch}{2.2}
\begin{tabularx}{\linewidth}{@{} >{\bfseries\color{navy}}p{0.32\linewidth} X @{}}
Two variables:           & \ans{Cups of coffee per week; GPA} \\
Variable types:          & \ans{Both quantitative} \\
Could this be random?    & \ans{Yes --- with only 30 students, this pattern could easily occur by chance} \\
One confounding variable: & \ans{Accept reasonable answers, e.g., hours of study (students who study more may drink more coffee \emph{and} have higher GPAs); sleep habits; socioeconomic background} \\
\end{tabularx}

\vspace{0.18in}

\textbf{Does this study prove causation?}\\[6pt]
\ans{No. This is an observational study, not a controlled experiment. Even if the
association is real, a confounding variable (such as amount of studying) could
explain it. To establish causation, we would need a randomized experiment in which
coffee consumption was randomly assigned.}

\vspace{0.2in}

\begin{tcolorbox}[colback=redbg, colframe=redacc, boxrule=0.8pt, arc=2mm,
    left=4mm, right=4mm, top=3mm, bottom=3mm]
\small\textbf{AP-Style Practice}\\[6pt]
A study found a strong positive association between the number of fire trucks sent to
a fire and the amount of property damage. Which of the following is the most plausible
explanation for this association?

\vspace{10pt}
\begin{enumerate}[label=\textbf{(\Alph*)}, leftmargin=2em, itemsep=6pt]
  \item Fire trucks cause property damage.
  \item Property damage attracts more fire trucks.
  \item \textbf{\textcolor{keyred}{The size of the fire is a confounding variable.}}
        \textcolor{keyred}{\textit{$\leftarrow$ correct}}
  \item The association is probably random because no experiment was conducted.
\end{enumerate}
\end{tcolorbox}

\vspace{0.1in}
\noindent\ans{Answer: C} \quad
\ans{Larger fires both require more trucks and cause more damage. The fire's size is
related to both variables, making it a confounding variable. The association between
trucks and damage is real but not causal. (D) is wrong --- a real association
can exist in observational data; the issue is causation, not randomness.}

\vspace{0.12in}
\begin{teachernote}
\textbf{Scoring (completeness):}
Award credit for any good-faith attempt. Common gaps:

\textbf{Confounding variable for coffee study:} Students who say ``genetics'' or
``motivation'' are acceptable --- they are variables related to both. Push back gently
on answers that only relate to one variable (e.g., ``liking the taste of coffee'').

\textbf{AP item:} Most common wrong answer is A (surface-level causal reading) or B
(reverse causation). If more than a third of the class misses this, revisit the
three-explanations framework (X causes Y / Y causes X / third variable causes both)
at the start of Lesson 2.2.

\textbf{D is a subtle distractor:} The association \emph{is} real --- fires do require more
trucks --- it's just that the relationship is confounded by fire size, not random.
Distinguishing ``random association'' from ``confounded association'' is a key conceptual
target for this lesson.
\end{teachernote}

\end{document}
